\documentclass[12pt]{article}

\usepackage[margin=1in]{geometry}
\usepackage{amsmath, amsfonts, amssymb}
\usepackage{enumitem}

\setlength{\parindent}{0pt}
\setlength{\parskip}{6pt}

\newtheorem{theorem}{Theorem}
\newtheorem{lemma}{Lemma}

\begin{document}

\section*{Axiomatic System}

\textbf{Axiom 1 (Infinite Adjunction).}
Beginning with a base set \(P_0\), whenever a set \(P_n\) has been constructed, there exists a fixed distinguished element
\[
\alpha_{n+1}\notin P_n
\]
which has not appeared at any preceding construction stage. Distinct stages use distinct elements. Hence the construction may continue indefinitely.

At this point the subscripts
\[
0,1,2,\ldots
\]
serve only as construction-stage labels. Their arithmetic properties have not yet been assumed or defined.

\textbf{Axiom 2 (Function Well-Definedness).}
A function \(f:X\to Y\) maps every input \(x\in X\) to exactly one output \(f(x)\in Y\). We denote evaluation by
\[
f(x)\to y.
\]

\textbf{Axiom 3 (Structural Induction).}
Let \(\mathcal P\) be a family generated from a base object \(P_0\) by repeated application of a function \(S\). If a property holds for \(P_0\), and whenever it holds for a constructed set \(P_n\) it also holds for \(S(P_n)\), then the property holds for every set generated from \(P_0\) by repeated applications of \(S\).

\section*{Definitions and Rules}

\subsection*{1.1 Base Set}

Let
\[
P_0=\varnothing.
\]

The empty set serves as the base representation of the symbol
\[
0.
\]

Thus
\[
E(P_0)=0.
\]

\subsection*{1.2 Canonical Finite Sets}

Using Axiom 1, define the canonical sets recursively.

Whenever \(P_n\) has already been constructed, define
\[
P_{n+1}=P_n\cup\{\alpha_{n+1}\},
\]
where
\[
\alpha_{n+1}\notin P_n.
\]

Consequently,
\[
P_0=\varnothing,
\]
\[
P_1=\{\alpha_1\},
\]
\[
P_2=\{\alpha_1,\alpha_2\},
\]
\[
P_3=\{\alpha_1,\alpha_2,\alpha_3\},
\]
and so on.

Let \(\mathcal P\) denote the smallest family containing \(P_0\) and closed under this construction.

Thus
\[
\mathcal P=\{P_0,P_1,P_2,P_3,\ldots\}.
\]

Because the element introduced at each stage is fixed, each construction stage has a unique output.

\subsection*{1.3 Output Symbols and the Cardinality Function}

We now assign one output symbol to each construction stage.

Define
\[
E:\mathcal P\to \{0,1,2,3,\ldots\}
\]
recursively by
\[
E(P_0)=0,
\]
and whenever \(P_{n+1}\) is the immediate successor of \(P_n\), let \(E(P_{n+1})\) be the next output symbol after \(E(P_n)\).

Accordingly,
\[
E(P_0)=0,
\]
\[
E(P_1)=1,
\]
\[
E(P_2)=2,
\]
\[
E(P_3)=3,
\]
and so on.

Hence, after these symbols have been generated,
\[
E(P_n)=n.
\]

The statement \(E(P_n)=n\) therefore does not use a previously defined arithmetic operation. It associates the symbol generated at stage \(n\) with the canonical set generated at the same stage.

Since each \(P_n\) contains exactly the elements introduced during the first \(n\) nonzero construction stages, \(E(P_n)\) is also the finite cardinality of \(P_n\). Therefore, once the symbols have been constructed, we may write
\[
E(P_n)=|P_n|.
\]

The domain of \(E\) in this construction is \(\mathcal P\), so \(E\) is never required to assign a finite natural-number output to an arbitrary infinite set.

\subsection*{1.4 Successor Function}

Define
\[
S:\mathcal P\to\mathcal P
\]
by
\[
S(P_n)=P_n\cup\{\alpha_{n+1}\}=P_{n+1}.
\]

Because \(\alpha_{n+1}\) is fixed at construction stage \(n+1\), \(S\) maps every canonical set to a unique canonical set and is therefore well-defined.

Applying \(S\) repeatedly gives
\[
P_0=\varnothing,
\]
\[
S(P_0)=P_1=\{\alpha_1\},
\]
\[
S(S(P_0))=P_2=\{\alpha_1,\alpha_2\},
\]
\[
S(S(S(P_0)))=P_3=\{\alpha_1,\alpha_2,\alpha_3\}.
\]

Evaluating with \(E\),
\[
E(P_0)=0,
\]
\[
E(S(P_0))=1,
\]
\[
E(S(S(P_0)))=2,
\]
\[
E(S(S(S(P_0))))=3.
\]

For any already constructed natural-number symbol \(k\), write
\[
S^k(P)
\]
for \(k\) repeated applications of \(S\), with
\[
S^0(P)=P.
\]

\subsection*{1.5 Natural Numbers Definition}

The natural numbers are defined to be exactly the output symbols produced by evaluating \(E\) on the canonical sets:
\[
\mathbb N
=
\{E(P)\mid P\in\mathcal P\}.
\]

Thus
\[
\boxed{
\mathbb N=\{0,1,2,3,\ldots\}.
}
\]

This definition does not require a statement of the form \(k\in\mathbb N\) before \(\mathbb N\) itself has been constructed.

\section*{Defining Addition}

We define addition using the canonical sets, the successor function \(S\), and the evaluation function \(E\).

For \(P_1,P_2\in\mathcal P\), define
\[
\boxed{
A(P_1,P_2)
=
E\left(S^{E(P_2)}(P_1)\right).
}
\]

Thus \(A(P_1,P_2)\) is a natural-number symbol, not a set.

We define the ordinary notation \(+\) by
\[
\boxed{
E(P_1)+E(P_2)
:=
A(P_1,P_2).
}
\]

Therefore,
\[
E(P_1)+E(P_2)
=
E\left(S^{E(P_2)}(P_1)\right).
\]

\subsection*{Addition by Zero}

Since
\[
E(P_0)=0,
\]
we have
\[
A(P_1,P_0)
=
E\left(S^{E(P_0)}(P_1)\right)
=
E(S^0(P_1))
=
E(P_1).
\]

Hence
\[
\boxed{
E(P_1)+0=E(P_1).
}
\]

Also,
\[
A(P_0,P_1)
=
E\left(S^{E(P_1)}(P_0)\right).
\]

By construction, applying \(S\) to \(P_0\) exactly \(E(P_1)\) times produces \(P_1\). Therefore
\[
A(P_0,P_1)=E(P_1),
\]
so
\[
\boxed{
0+E(P_1)=E(P_1).
}
\]

\subsection*{Successor Compatibility}

Suppose
\[
A(P_1,P_2)=E(P_r).
\]

Then
\[
S^{E(P_2)}(P_1)=P_r.
\]

Since \(S(P_2)\) contains exactly one more constructed element than \(P_2\),
\[
S^{E(S(P_2))}(P_1)
=
S\left(S^{E(P_2)}(P_1)\right).
\]

Therefore
\[
A(P_1,S(P_2))
=
E\left(
S\left(S^{E(P_2)}(P_1)\right)
\right)
=
E(S(P_r)).
\]

Likewise,
\[
S^{E(P_1)}(S(P_2))
=
S\left(S^{E(P_1)}(P_2)\right).
\]

Thus taking one successor step in either input produces one successor step in the resulting canonical set.

\subsection*{Commutativity of Addition}

\begin{theorem}
For all \(P_1,P_2\in\mathcal P\),
\[
A(P_1,P_2)=A(P_2,P_1).
\]
Consequently,
\[
E(P_1)+E(P_2)=E(P_2)+E(P_1).
\]
\end{theorem}

\begin{proof}
Fix \(P_1\). We use structural induction on \(P_2\).

For the base case,
\[
P_2=P_0.
\]

Then
\[
A(P_1,P_0)=E(P_1).
\]

Also,
\[
A(P_0,P_1)
=
E\left(S^{E(P_1)}(P_0)\right)
=
E(P_1).
\]

Hence
\[
A(P_1,P_0)=A(P_0,P_1).
\]

Now assume
\[
A(P_1,P_2)=A(P_2,P_1).
\]

Let
\[
P_2'=S(P_2).
\]

By successor compatibility,
\[
A(P_1,P_2')
\]
is the output symbol immediately following
\[
A(P_1,P_2).
\]

Likewise,
\[
A(P_2',P_1)
\]
is the output symbol immediately following
\[
A(P_2,P_1).
\]

By the induction hypothesis,
\[
A(P_1,P_2)=A(P_2,P_1),
\]
so their immediate successor symbols are equal. Therefore
\[
A(P_1,P_2')=A(P_2',P_1).
\]

By structural induction,
\[
A(P_1,P_2)=A(P_2,P_1)
\]
for every \(P_2\in\mathcal P\).
\end{proof}

Therefore it is valid to write
\[
\boxed{
E(P_1)+E(P_2)
=
E(P_2)+E(P_1).
}
\]

\subsection*{Example: \(1+2\)}

Let
\[
P_0=\varnothing,
\]
\[
P_1=S(P_0),
\]
\[
P_2=S(S(P_0)).
\]

Then
\[
E(P_1)=1,
\qquad
E(P_2)=2.
\]

Evaluate
\[
A(P_1,P_2)
=
E\left(S^{E(P_2)}(P_1)\right).
\]

Since
\[
E(P_2)=2,
\]
we obtain
\[
A(P_1,P_2)
=
E(S^2(P_1)).
\]

But
\[
S^2(P_1)
=
S(S(P_1))
=
P_3.
\]

Thus
\[
A(P_1,P_2)
=
E(P_3)
=
3.
\]

Therefore
\[
1+2=3.
\]

By the commutativity theorem,
\[
A(P_2,P_1)=A(P_1,P_2)=3,
\]
so
\[
2+1=3.
\]

Hence
\[
\boxed{
1+2=2+1=3.
}
\]

\section*{Subtraction on Natural Numbers via Element Removal}

We now define subtraction when the result remains nonnegative.

\subsection*{1. Element Removal Function}

Define the removal function \(N\) on every nonempty canonical set by
\[
\boxed{
N(P_n)=P_{n-1}
\qquad\text{for }P_n\neq P_0.
}
\]

Equivalently, \(N\) removes the distinguished element introduced most recently in the canonical construction.

For example,
\[
P_3=\{\alpha_1,\alpha_2,\alpha_3\}.
\]

Then
\[
N(P_3)
=
\{\alpha_1,\alpha_2\}
=
P_2.
\]

Similarly,
\[
N(P_2)=P_1
\]
and
\[
N(P_1)=P_0.
\]

In particular,
\[
N(P_1)=P_0=\varnothing,
\]
and therefore
\[
E(N(P_1))=0.
\]

The rule specifies a unique element to remove, so \(N\) is well-defined.

\subsection*{2. Iterated Removal}

For any natural-number symbol \(k\), define
\[
N^0(P)=P,
\]
and, whenever the intermediate set is nonempty,
\[
N^k(P)
=
N\left(N^{k-1}(P)\right)
\]
for nonzero \(k\).

For example,
\[
N^2(P_3)
=
N(N(P_3))
=
N(P_2)
=
P_1.
\]

\subsection*{3. Ordering of Natural Numbers}

Before restricting subtraction, define the order relation \(\leq\).

For canonical sets \(P_1,P_2\), define
\[
\boxed{
E(P_2)\leq E(P_1)
}
\]
to mean that \(P_1\) can be obtained from \(P_2\) by zero or more applications of \(S\).

Equivalently,
\[
E(P_2)\leq E(P_1)
\]
when there exists a natural-number symbol \(k\) such that
\[
S^k(P_2)=P_1.
\]

Thus, for example,
\[
1\leq3
\]
because
\[
S^2(P_1)=P_3.
\]

\subsection*{4. Definition of Natural Subtraction}

Let
\[
D(P_1,P_2)
\]
be defined whenever
\[
E(P_2)\leq E(P_1).
\]

Define
\[
\boxed{
D(P_1,P_2)
=
E\left(
N^{E(P_2)}(P_1)
\right).
}
\]

We then define
\[
\boxed{
E(P_1)-E(P_2)
:=
D(P_1,P_2)
}
\]
whenever
\[
E(P_2)\leq E(P_1).
\]

The restriction guarantees that \(N\) is never required to remove an element from \(P_0\).

\subsection*{5. Example Evaluation}

Let
\[
P_3=\{\alpha_1,\alpha_2,\alpha_3\}
\]
and
\[
P_1=\{\alpha_1\}.
\]

Then
\[
E(P_3)=3
\]
and
\[
E(P_1)=1.
\]

Evaluate
\[
D(P_3,P_1)
=
E\left(N^{E(P_1)}(P_3)\right).
\]

Since
\[
E(P_1)=1,
\]
we have
\[
D(P_3,P_1)
=
E(N(P_3)).
\]

By the definition of \(N\),
\[
N(P_3)=P_2.
\]

Therefore
\[
D(P_3,P_1)
=
E(P_2)
=
2.
\]

Thus
\[
\boxed{
3-1=2.
}
\]

\section*{Constructing Negative Numbers}

Natural-number subtraction as defined above does not yet assign an output to expressions such as
\[
1-2.
\]

We now construct signed integers using ordered pairs of canonical sets.

\subsection*{1. Ordered-Pair Representation}

Take an ordered pair
\[
(P_x,P_y),
\]
where
\[
P_x,P_y\in\mathcal P.
\]

The first component records a positive contribution, and the second component records a negative contribution.

At this stage, however, the symbol ``negative'' describes the role of the second component only. A negative-number symbol such as \(-1\) will be introduced only after the corresponding ordered-pair object has been constructed.

\subsection*{2. Reduction Function}

Define
\[
R[(P_x,P_y)]
\]
recursively by
\[
\boxed{
R[(P_x,P_y)]
=
\begin{cases}
R[(N(P_x),N(P_y))],
&
E(P_x)>0\text{ and }E(P_y)>0,
\\[6pt]
(P_x,P_y),
&
E(P_x)=0\text{ or }E(P_y)=0.
\end{cases}
}
\]

Here
\[
E(P)>0
\]
means simply
\[
P\neq P_0.
\]

Every recursive step removes one element from both components.

Since \(P_x\) and \(P_y\) are finite canonical sets, after finitely many applications of \(N\), at least one component must become \(P_0\).

Therefore \(R\) always terminates.

Most importantly,
\[
R
\]
always returns an ordered pair. Its output type does not change during the definition.

\subsection*{3. Reduced Signed Pairs}

A pair
\[
(P_x,P_y)
\]
is called \emph{reduced} if at least one component equals \(P_0\).

Thus every reduced pair has exactly one of the following forms:
\[
(P_n,P_0),
\]
\[
(P_0,P_n),
\]
or
\[
(P_0,P_0).
\]

The function \(R\) maps every ordered pair of canonical sets to exactly one reduced pair.

\subsection*{4. Signed Integer Notation}

We now introduce notation for these already-constructed reduced pairs.

For a nonzero natural-number symbol \(n\), define the notation
\[
(P_n,P_0)\equiv n.
\]

Define
\[
(P_0,P_0)\equiv0.
\]

Finally, for a nonzero natural-number symbol \(n\), introduce the new notation
\[
\boxed{
(P_0,P_n)\equiv -n.
}
\]

Thus the symbol
\[
-n
\]
is not used to construct the negative integer.

Instead, the ordered pair
\[
(P_0,P_n)
\]
is constructed first, and \(-n\) is then assigned as notation for that object.

The signed integers are therefore the reduced pairs
\[
\boxed{
\mathbb Z
=
\{
R[(P_x,P_y)]
\mid
P_x,P_y\in\mathcal P
\}.
}
\]

Using the notation above, these objects are written
\[
\mathbb Z
=
\{\ldots,-3,-2,-1,0,1,2,3,\ldots\}.
\]

\subsection*{5. Canonical Representation}

Because \(R\) removes one element from each side simultaneously, every pair has a unique reduced representative.

For example,
\[
R[(P_2,P_1)]
=
R[(N(P_2),N(P_1))]
=
R[(P_1,P_0)]
=
(P_1,P_0).
\]

Thus
\[
(P_2,P_1)
\]
and
\[
(P_1,P_0)
\]
reduce to the same signed integer.

More generally, two ordered pairs
\[
(P_a,P_b)
\]
and
\[
(P_c,P_d)
\]
represent the same signed integer precisely when
\[
\boxed{
R[(P_a,P_b)]
=
R[(P_c,P_d)].
}
\]

Hence signed integers have unique reduced representatives even though unreduced intermediate pairs may differ.

\subsection*{6. Examples}

\textbf{Example 1: \(1-2\).}

Take
\[
P_1=\{\alpha_1\}
\]
and
\[
P_2=\{\alpha_1,\alpha_2\}.
\]

Evaluate
\[
R[(P_1,P_2)].
\]

Since both components are nonempty,
\[
R[(P_1,P_2)]
=
R[(N(P_1),N(P_2))].
\]

Using
\[
N(P_1)=P_0
\]
and
\[
N(P_2)=P_1,
\]
we obtain
\[
R[(P_1,P_2)]
=
R[(P_0,P_1)].
\]

Because the first component is now \(P_0\),
\[
R[(P_0,P_1)]
=
(P_0,P_1).
\]

By the signed notation introduced above,
\[
(P_0,P_1)\equiv-1.
\]

Therefore
\[
\boxed{
1-2=-1.
}
\]

\textbf{Example 2: \(2-1\).}

Evaluate
\[
R[(P_2,P_1)].
\]

Since both components are nonempty,
\[
R[(P_2,P_1)]
=
R[(N(P_2),N(P_1))].
\]

Thus
\[
R[(P_2,P_1)]
=
R[(P_1,P_0)]
=
(P_1,P_0).
\]

By signed notation,
\[
(P_1,P_0)\equiv1.
\]

Therefore
\[
\boxed{
2-1=1.
}
\]

\textbf{Example 3: \(2-2\).}

Evaluate
\[
R[(P_2,P_2)].
\]

First,
\[
R[(P_2,P_2)]
=
R[(N(P_2),N(P_2))]
=
R[(P_1,P_1)].
\]

Both components remain nonempty, so
\[
R[(P_1,P_1)]
=
R[(N(P_1),N(P_1))]
=
R[(P_0,P_0)].
\]

The pair is now reduced, giving
\[
R[(P_0,P_0)]
=
(P_0,P_0).
\]

Since
\[
(P_0,P_0)\equiv0,
\]
we obtain
\[
\boxed{
2-2=0.
}
\]

\section*{Addition and Subtraction of Signed Pairs}

We now define arithmetic directly on reduced signed-pair representations.

Let
\[
(P_a,P_b)
\]
and
\[
(P_c,P_d)
\]
be signed-pair representations.

The positive contributions are represented by
\[
P_a
\quad\text{and}\quad
P_c,
\]
while the negative contributions are represented by
\[
P_b
\quad\text{and}\quad
P_d.
\]

\subsection*{Signed Addition}

The combined positive count is
\[
A(P_a,P_c).
\]

The corresponding canonical set is
\[
S^{A(P_a,P_c)}(P_0).
\]

Likewise, the combined negative count is
\[
A(P_b,P_d),
\]
with corresponding canonical set
\[
S^{A(P_b,P_d)}(P_0).
\]

Therefore define
\[
\boxed{
\operatorname{Add}
\left(
(P_a,P_b),(P_c,P_d)
\right)
=
R\left[
\left(
S^{A(P_a,P_c)}(P_0),
S^{A(P_b,P_d)}(P_0)
\right)
\right].
}
\]

The result is a reduced signed pair.

\subsection*{Signed Additive Inverse}

For a signed pair
\[
(P_c,P_d),
\]
define its additive inverse by exchanging its two components:
\[
\boxed{
-(P_c,P_d):=(P_d,P_c).
}
\]

This is valid because the component that formerly contributed positively now contributes negatively, and the component that formerly contributed negatively now contributes positively.

For example,
\[
-(P_2,P_0)
=
(P_0,P_2),
\]
so the additive inverse of \(2\) is represented by
\[
(P_0,P_2)\equiv-2.
\]

Similarly,
\[
-(P_0,P_2)
=
(P_2,P_0),
\]
so the additive inverse of \(-2\) is \(2\).

\subsection*{Signed Subtraction}

Subtraction is addition of the additive inverse.

Thus
\[
(P_a,P_b)-(P_c,P_d)
\]
is defined as
\[
(P_a,P_b)+(P_d,P_c).
\]

Therefore
\[
\boxed{
\operatorname{Subtract}
\left(
(P_a,P_b),(P_c,P_d)
\right)
=
R\left[
\left(
S^{A(P_a,P_d)}(P_0),
S^{A(P_b,P_c)}(P_0)
\right)
\right].
}
\]

\subsection*{1. Adding Same Signs: \((-2)+(-1)\)}

The inputs are
\[
(P_0,P_2)
\]
and
\[
(P_0,P_1).
\]

The new positive count is
\[
A(P_0,P_0)=0.
\]

Therefore
\[
P_{\mathrm{pos}}
=
S^0(P_0)
=
P_0.
\]

The new negative count is
\[
A(P_2,P_1)=3.
\]

Therefore
\[
P_{\mathrm{neg}}
=
S^3(P_0)
=
P_3.
\]

Hence
\[
\operatorname{Add}
\left(
(P_0,P_2),(P_0,P_1)
\right)
=
R[(P_0,P_3)].
\]

Since the first component is already \(P_0\),
\[
R[(P_0,P_3)]
=
(P_0,P_3).
\]

By signed notation,
\[
(P_0,P_3)\equiv-3.
\]

Therefore
\[
\boxed{
(-2)+(-1)=-3.
}
\]

\subsection*{2. Adding Opposite Signs: \((-2)+1\)}

The inputs are
\[
(P_0,P_2)
\]
and
\[
(P_1,P_0).
\]

The new positive count is
\[
A(P_0,P_1)=1,
\]
so
\[
P_{\mathrm{pos}}=P_1.
\]

The new negative count is
\[
A(P_2,P_0)=2,
\]
so
\[
P_{\mathrm{neg}}=P_2.
\]

Thus
\[
\operatorname{Add}
\left(
(P_0,P_2),(P_1,P_0)
\right)
=
R[(P_1,P_2)].
\]

Applying \(N\) once to both components gives
\[
R[(P_1,P_2)]
=
R[(P_0,P_1)]
=
(P_0,P_1).
\]

Since
\[
(P_0,P_1)\equiv-1,
\]
we obtain
\[
\boxed{
(-2)+1=-1.
}
\]

\subsection*{3. Subtracting Same Signs: \(1-2\)}

The inputs are
\[
(P_1,P_0)
\]
and
\[
(P_2,P_0).
\]

Using the subtraction definition,
\[
A(P_1,P_0)=1,
\]
so
\[
P_{\mathrm{pos}}=P_1.
\]

Also,
\[
A(P_0,P_2)=2,
\]
so
\[
P_{\mathrm{neg}}=P_2.
\]

Thus
\[
\operatorname{Subtract}
\left(
(P_1,P_0),(P_2,P_0)
\right)
=
R[(P_1,P_2)].
\]

Reduction gives
\[
R[(P_1,P_2)]
=
R[(P_0,P_1)]
=
(P_0,P_1).
\]

Since
\[
(P_0,P_1)\equiv-1,
\]
we obtain
\[
\boxed{
1-2=-1.
}
\]

\subsection*{4. Subtracting Opposite Signs: \(1-(-2)\)}

The inputs are
\[
(P_1,P_0)
\]
and
\[
(P_0,P_2).
\]

The positive count is
\[
A(P_1,P_2)=3,
\]
so
\[
P_{\mathrm{pos}}=P_3.
\]

The negative count is
\[
A(P_0,P_0)=0,
\]
so
\[
P_{\mathrm{neg}}=P_0.
\]

Therefore
\[
\operatorname{Subtract}
\left(
(P_1,P_0),(P_0,P_2)
\right)
=
R[(P_3,P_0)].
\]

Since the second component is already \(P_0\),
\[
R[(P_3,P_0)]
=
(P_3,P_0).
\]

By signed notation,
\[
(P_3,P_0)\equiv3.
\]

Hence
\[
\boxed{
1-(-2)=3.
}
\]

\section*{Proof by Fully Expanded Substitution}

We now expand the preceding signed arithmetic computations back to the base set \(P_0\) using only
\[
S(P),\qquad
E(P),\qquad
A(P_1,P_2),\qquad
N(P),\qquad
R[(P_x,P_y)].
\]

\subsection*{1. Adding Same Signs: \((-2)+(-1)=-3\)}

The base representations are
\[
\text{Pair 1}
=
(P_0,S(S(P_0)))
\]
and
\[
\text{Pair 2}
=
(P_0,S(P_0)).
\]

The expanded addition formula is
\[
\operatorname{Add}
=
R\left[
\left(
S^{A(P_0,P_0)}(P_0),
S^{A(S(S(P_0)),S(P_0))}(P_0)
\right)
\right].
\]

First evaluate the positive component:
\[
A(P_0,P_0)
=
E\left(
S^{E(P_0)}(P_0)
\right).
\]

Since
\[
E(P_0)=0,
\]
we obtain
\[
A(P_0,P_0)
=
E(S^0(P_0))
=
E(P_0)
=
0.
\]

Thus
\[
k_{\mathrm{pos}}=0.
\]

Now evaluate the negative component:
\[
A(S(S(P_0)),S(P_0))
=
E\left(
S^{E(S(P_0))}(S(S(P_0)))
\right).
\]

Since
\[
E(S(P_0))=1,
\]
we obtain
\[
A(S(S(P_0)),S(P_0))
=
E\left(
S^1(S(S(P_0)))
\right).
\]

Therefore
\[
A(S(S(P_0)),S(P_0))
=
E(S(S(S(P_0))))
=
3.
\]

Thus
\[
k_{\mathrm{neg}}=3.
\]

Construct the new sets:
\[
P_x
=
S^0(P_0)
=
P_0,
\]
and
\[
P_y
=
S^3(P_0)
=
S(S(S(P_0))).
\]

Therefore
\[
\operatorname{Add}
=
R[(P_0,S(S(S(P_0))))].
\]

Since the first component is \(P_0\), reduction stops immediately:
\[
R[(P_0,S(S(S(P_0))))]
=
(P_0,S(S(S(P_0)))).
\]

But
\[
S(S(S(P_0)))=P_3.
\]

Hence the reduced pair is
\[
(P_0,P_3).
\]

By signed notation,
\[
(P_0,P_3)\equiv-3.
\]

Therefore
\[
\boxed{
(-2)+(-1)=-3.
}
\]

\subsection*{2. Adding Opposite Signs: \((-2)+1=-1\)}

The base representations are
\[
\text{Pair 1}
=
(P_0,S(S(P_0)))
\]
and
\[
\text{Pair 2}
=
(S(P_0),P_0).
\]

The expanded addition formula is
\[
\operatorname{Add}
=
R\left[
\left(
S^{A(P_0,S(P_0))}(P_0),
S^{A(S(S(P_0)),P_0)}(P_0)
\right)
\right].
\]

Evaluate the positive component:
\[
A(P_0,S(P_0))
=
E\left(
S^{E(S(P_0))}(P_0)
\right).
\]

Since
\[
E(S(P_0))=1,
\]
we have
\[
A(P_0,S(P_0))
=
E(S^1(P_0))
=
E(S(P_0))
=
1.
\]

Thus
\[
k_{\mathrm{pos}}=1.
\]

Now evaluate the negative component:
\[
A(S(S(P_0)),P_0)
=
E\left(
S^{E(P_0)}(S(S(P_0)))
\right).
\]

Since
\[
E(P_0)=0,
\]
we obtain
\[
A(S(S(P_0)),P_0)
=
E(S^0(S(S(P_0))))
=
E(S(S(P_0)))
=
2.
\]

Thus
\[
k_{\mathrm{neg}}=2.
\]

Construct the sets:
\[
P_x
=
S^1(P_0)
=
S(P_0),
\]
and
\[
P_y
=
S^2(P_0)
=
S(S(P_0)).
\]

Therefore
\[
\operatorname{Add}
=
R[(S(P_0),S(S(P_0)))].
\]

Because both sets are nonempty,
\[
R[(S(P_0),S(S(P_0)))]
=
R[(N(S(P_0)),N(S(S(P_0))))].
\]

By the definition of \(N\),
\[
N(S(P_0))=P_0
\]
and
\[
N(S(S(P_0)))=S(P_0).
\]

Therefore
\[
R[(S(P_0),S(S(P_0)))]
=
R[(P_0,S(P_0))].
\]

The first component is now \(P_0\), so
\[
R[(P_0,S(P_0))]
=
(P_0,S(P_0)).
\]

Since
\[
S(P_0)=P_1,
\]
the reduced pair is
\[
(P_0,P_1).
\]

By signed notation,
\[
(P_0,P_1)\equiv-1.
\]

Therefore
\[
\boxed{
(-2)+1=-1.
}
\]

\subsection*{3. Subtracting Same Signs: \(1-2=-1\)}

The base representations are
\[
\text{Pair 1}
=
(S(P_0),P_0)
\]
and
\[
\text{Pair 2}
=
(S(S(P_0)),P_0).
\]

The expanded subtraction formula is
\[
\operatorname{Subtract}
=
R\left[
\left(
S^{A(S(P_0),P_0)}(P_0),
S^{A(P_0,S(S(P_0)))}(P_0)
\right)
\right].
\]

Evaluate the positive component:
\[
A(S(P_0),P_0)
=
E\left(
S^{E(P_0)}(S(P_0))
\right).
\]

Since
\[
E(P_0)=0,
\]
we obtain
\[
A(S(P_0),P_0)
=
E(S^0(S(P_0)))
=
E(S(P_0))
=
1.
\]

Thus
\[
k_{\mathrm{pos}}=1.
\]

Evaluate the negative component:
\[
A(P_0,S(S(P_0)))
=
E\left(
S^{E(S(S(P_0)))}(P_0)
\right).
\]

Since
\[
E(S(S(P_0)))=2,
\]
we obtain
\[
A(P_0,S(S(P_0)))
=
E(S^2(P_0))
=
E(S(S(P_0)))
=
2.
\]

Thus
\[
k_{\mathrm{neg}}=2.
\]

Construct the sets:
\[
P_x
=
S^1(P_0)
=
S(P_0),
\]
and
\[
P_y
=
S^2(P_0)
=
S(S(P_0)).
\]

Therefore
\[
\operatorname{Subtract}
=
R[(S(P_0),S(S(P_0)))].
\]

Reduction gives
\[
R[(S(P_0),S(S(P_0)))]
=
R[(P_0,S(P_0))]
=
(P_0,S(P_0)).
\]

Since
\[
S(P_0)=P_1,
\]
we obtain
\[
(P_0,S(P_0))
=
(P_0,P_1).
\]

By signed notation,
\[
(P_0,P_1)\equiv-1.
\]

Therefore
\[
\boxed{
1-2=-1.
}
\]

\subsection*{4. Subtracting Opposite Signs: \(1-(-2)=3\)}

The base representations are
\[
\text{Pair 1}
=
(S(P_0),P_0)
\]
and
\[
\text{Pair 2}
=
(P_0,S(S(P_0))).
\]

The expanded subtraction formula is
\[
\operatorname{Subtract}
=
R\left[
\left(
S^{A(S(P_0),S(S(P_0)))}(P_0),
S^{A(P_0,P_0)}(P_0)
\right)
\right].
\]

Evaluate the positive component:
\[
A(S(P_0),S(S(P_0)))
=
E\left(
S^{E(S(S(P_0)))}(S(P_0))
\right).
\]

Since
\[
E(S(S(P_0)))=2,
\]
we obtain
\[
A(S(P_0),S(S(P_0)))
=
E(S^2(S(P_0))).
\]

Therefore
\[
A(S(P_0),S(S(P_0)))
=
E(S(S(S(P_0))))
=
3.
\]

Thus
\[
k_{\mathrm{pos}}=3.
\]

Evaluate the negative component:
\[
A(P_0,P_0)
=
E\left(
S^{E(P_0)}(P_0)
\right).
\]

Since
\[
E(P_0)=0,
\]
we obtain
\[
A(P_0,P_0)
=
E(P_0)
=
0.
\]

Thus
\[
k_{\mathrm{neg}}=0.
\]

Construct the sets:
\[
P_x
=
S^3(P_0)
=
S(S(S(P_0))),
\]
and
\[
P_y
=
S^0(P_0)
=
P_0.
\]

Therefore
\[
\operatorname{Subtract}
=
R[(S(S(S(P_0))),P_0)].
\]

Since the second component is already \(P_0\),
\[
R[(S(S(S(P_0))),P_0)]
=
(S(S(S(P_0))),P_0).
\]

Since
\[
S(S(S(P_0)))=P_3,
\]
the reduced pair is
\[
(P_3,P_0).
\]

By signed notation,
\[
(P_3,P_0)\equiv3.
\]

Therefore
\[
\boxed{
1-(-2)=3.
}
\]

\section*{Conclusion}

Beginning with
\[
P_0=\varnothing,
\]
the canonical sets are generated by the well-defined successor function
\[
S(P_n)=P_{n+1}.
\]

The evaluation function \(E\) associates each canonical construction stage with exactly one natural-number symbol:
\[
E(P_n)=n.
\]

This produces
\[
\mathbb N=\{0,1,2,3,\ldots\}.
\]

Addition is then defined by
\[
\boxed{
A(P_1,P_2)
=
E\left(
S^{E(P_2)}(P_1)
\right),
}
\]
with
\[
E(P_1)+E(P_2)
:=
A(P_1,P_2).
\]

Commutativity follows by structural induction:
\[
\boxed{
E(P_1)+E(P_2)
=
E(P_2)+E(P_1).
}
\]

For nonnegative results, subtraction is defined by the deterministic removal function
\[
N(P_n)=P_{n-1}
\]
and
\[
\boxed{
D(P_1,P_2)
=
E\left(
N^{E(P_2)}(P_1)
\right).
}
\]

Signed integers are then constructed as reduced ordered pairs of canonical sets. The reduction function
\[
\boxed{
R[(P_x,P_y)]
=
\begin{cases}
R[(N(P_x),N(P_y))],
&
E(P_x)>0\text{ and }E(P_y)>0,
\\[6pt]
(P_x,P_y),
&
E(P_x)=0\text{ or }E(P_y)=0
\end{cases}
}
\]
always returns a reduced ordered pair.

Only after these ordered-pair objects have been constructed is signed notation introduced:
\[
(P_n,P_0)\equiv n,
\]
\[
(P_0,P_0)\equiv0,
\]
and
\[
(P_0,P_n)\equiv-n.
\]

Signed addition is defined by
\[
\boxed{
\operatorname{Add}
\left(
(P_a,P_b),(P_c,P_d)
\right)
=
R\left[
\left(
S^{A(P_a,P_c)}(P_0),
S^{A(P_b,P_d)}(P_0)
\right)
\right],
}
\]
while signed subtraction is defined by
\[
\boxed{
\operatorname{Subtract}
\left(
(P_a,P_b),(P_c,P_d)
\right)
=
R\left[
\left(
S^{A(P_a,P_d)}(P_0),
S^{A(P_b,P_c)}(P_0)
\right)
\right].
}
\]

The explicit evaluations establish
\[
1+2=3,
\]
\[
3-1=2,
\]
\[
1-2=-1,
\]
\[
2-1=1,
\]
\[
2-2=0,
\]
\[
(-2)+(-1)=-3,
\]
\[
(-2)+1=-1,
\]
and
\[
1-(-2)=3.
\]

Thus the natural-number and signed-integer operations used above are obtained from the base set \(P_0\), deterministic adjunction and removal, evaluation of canonical finite sets, and reduction of ordered pairs.

\end{document}